\documentclass[12pt,a4paper]{article}

% --- Pakiety podstawowe ---
\usepackage[polish]{babel}
\usepackage[T1]{fontenc}
\usepackage[utf8]{inputenc}
\usepackage{amsmath}
\usepackage{amsfonts}
\usepackage{amssymb}
\usepackage{graphicx}
\usepackage{geometry}
\usepackage{xcolor}
\usepackage{titlesec}
\usepackage{fancyhdr}
\usepackage{longtable}
\usepackage{booktabs}
\usepackage{array}
\usepackage{hyperref}
\usepackage{enumitem}
\usepackage{calc}
% Pandoc-generated macro
\providecommand{\tightlist}{\setlength{\itemsep}{0pt}\setlength{\parskip}{0pt}}

% --- Konfiguracja strony ---
\geometry{margin=2.5cm}
\setlength{\parindent}{0pt}
\setlength{\parskip}{1em}

% --- Kolorystyka ---
\definecolor{NavyBlue}{RGB}{31, 56, 100}
\definecolor{MidBlue}{RGB}{46, 117, 182}
\definecolor{LightBlue}{RGB}{214, 228, 240}
\definecolor{Gray}{RGB}{242, 242, 242}

% --- Stylizacja nagłówków ---
\titleformat{\section}
  {\color{NavyBlue}\normalfont\Large\bfseries}{\thesection}{1em}{}[{\titlerule[1.5pt]}]

\titleformat{\subsection}
  {\color{MidBlue}\normalfont\large\bfseries}{\thesubsection}{1em}{}

% --- Nagłówek i stopka ---
\pagestyle{fancy}
\fancyhf{}
\fancyhead[L]{\small \color{gray} Pakiet Dokumentacji Cyberbezpieczeństwa v1.0}
\fancyhead[R]{\small \color{gray} [NAZWA FIRMY]}
\fancyfoot[C]{\thepage}
\fancyfoot[R]{\small \color{gray} POUFNE}

% --- Ustawienia tabel ---
\setlength{\LTpre}{4pt}
\setlength{\LTpost}{4pt}
\newcolumntype{L}[1]{>{\raggedright\arraybackslash}p{#1}}
\renewcommand{\arraystretch}{1.5}

\begin{document}

% --- STRONA TYTUŁOWA ---
\begin{titlepage}
    \centering
    \vspace*{2cm}
    {\huge\bfseries\color{NavyBlue} [NAZWA FIRMY KSIĘGOWEJ] \par}
    \vspace{1cm}
    {\Large\bfseries\color{MidBlue} PAKIET POLITYK I PROCEDUR\\CYBERBEZPIECZEŃSTWA \par}
    \vspace{2cm}
    
    \begin{table}[h]
        \centering
        \begin{tabular}{|>{\color{NavyBlue}\bfseries}l|l|}
            \hline
            Właściciel dokumentu: & Administrator / Właściciel firmy \\ \hline
            Status: & SZABLON -- do wypełnienia \\ \hline
            Data utworzenia: & [DATA] \\ \hline
            Następny przegląd: & [DATA + 12 MIESIĘCY] \\ \hline
            Wersja: & 1.0 \\ \hline
            Klasyfikacja: & POUFNE -- wewnętrzne \\ \hline
        \end{tabular}
    \end{table}

    \vfill
    \begin{quote}
        \centering
        \textit{\color{red} Dokument zawiera informacje poufne. Dystrybucja wyłącznie dla upoważnionych pracowników.}
    \end{quote}
    \vspace*{1cm}
\end{titlepage}

\newpage




\section{SPIS TREŚCI}\label{spis-treux15bci}

\begin{longtable}[]{@{}>{\raggedright\arraybackslash}p{\dimexpr\linewidth/1000*270-2\tabcolsep\relax}>{\raggedright\arraybackslash}p{\dimexpr\linewidth/1000*430-2\tabcolsep\relax}>{\raggedright\arraybackslash}p{\dimexpr\linewidth/1000*300-2\tabcolsep\relax}@{}}
\toprule\noalign{}
\endhead
\bottomrule\noalign{}
\endlastfoot
1 & Polityka przechowywania i zarządzania danymi (RODO) & POL-001 \\
2 & Procedura weryfikacji certyfikacji sprzętu & PROC-002 \\
3 & Procedura weryfikacji certyfikacji oprogramowania & PROC-003 \\
4 & Procedura zabezpieczania infrastruktury & PROC-004 \\
5 & Polityka dostępu i monitorowania użytkowników & POL-005 \\
6 & Procedura onboardingu nowych użytkowników & PROC-006 \\
7 & Procedura offboardingu odchodzących użytkowników & PROC-007 \\
8 & Procedura szkoleń i weryfikacji wiedzy & PROC-008 \\
9 & Polityka monitorowania i audytowania infrastruktury & POL-009 \\
10 & Playbook sytuacji nietypowych (Incident Response) & PLAY-010 \\
11 & Plan odtwarzania po katastrofie (Disaster Recovery) & DRP-011 \\
12 & Plan ciągłego doskonalenia cyberbezpieczeństwa & PLAN-012 \\
\end{longtable}

\phantomsection\label{doc1}
POL-001

\section{POLITYKA PRZECHOWYWANIA I ZARZĄDZANIA DANYMI
(RODO)}\label{polityka-przechowywania-i-zarzux105dzania-danymi-rodo}

\begin{longtable}[]{@{}>{\raggedright\arraybackslash}p{\dimexpr\linewidth/1000*300-2\tabcolsep\relax}>{\raggedright\arraybackslash}p{\dimexpr\linewidth/1000*700-2\tabcolsep\relax}@{}}
\toprule\noalign{}
\endhead
\bottomrule\noalign{}
\endlastfoot
Wersja: & 1.0 \\
Data wejścia w życie: & {[}DATA{]} \\
Właściciel: & Właściciel firmy / Inspektor Ochrony Danych (IOD) \\
Dotyczy: & Wszyscy pracownicy, współpracownicy, podwykonawcy \\
Przegląd: & Co 12 miesięcy lub po istotnej zmianie prawnej \\
\end{longtable}

\subsection{1. Cel i zakres}\label{cel-i-zakres}

Niniejsza polityka określa zasady gromadzenia, przetwarzania,
przechowywania i usuwania danych osobowych oraz danych finansowych
klientów w {[}NAZWA FIRMY{]}, zgodnie z Rozporządzeniem RODO (UE)
2016/679 oraz polskim prawem krajowym.

Polityka obejmuje: dane osobowe klientów i pracowników, dane finansowe i
księgowe, dokumenty elektroniczne i papierowe, systemy informatyczne i
nośniki danych.

\subsection{2. Podstawy prawne przetwarzania
danych}\label{podstawy-prawne-przetwarzania-danych}

\begin{longtable}[]{@{}>{\raggedright\arraybackslash}p{\dimexpr\linewidth/1000*270-2\tabcolsep\relax}>{\raggedright\arraybackslash}p{\dimexpr\linewidth/1000*430-2\tabcolsep\relax}>{\raggedright\arraybackslash}p{\dimexpr\linewidth/1000*300-2\tabcolsep\relax}@{}}
\toprule\noalign{}
Kategoria danych & Cel przetwarzania & Podstawa prawna (RODO) \\
\midrule\noalign{}
\endhead
\bottomrule\noalign{}
\endlastfoot
Dane klientów (osoby fizyczne) & Wykonanie usług księgowych & Art. 6
ust. 1 lit. b) i c) \\
Dane pracowników & Realizacja umowy o pracę & Art. 6 ust. 1 lit. b) i
c) \\
Dane do celów podatkowych & Obowiązek prawny (Ordynacja podatkowa) &
Art. 6 ust. 1 lit. c) \\
Dane marketingowe & Zgoda osoby & Art. 6 ust. 1 lit. a) \\
\end{longtable}

\subsection{3. Okresy przechowywania
danych}\label{okresy-przechowywania-danych}

\begin{longtable}[]{@{}>{\raggedright\arraybackslash}p{\dimexpr\linewidth/1000*270-2\tabcolsep\relax}>{\raggedright\arraybackslash}p{\dimexpr\linewidth/1000*430-2\tabcolsep\relax}>{\raggedright\arraybackslash}p{\dimexpr\linewidth/1000*300-2\tabcolsep\relax}@{}}
\toprule\noalign{}
Rodzaj dokumentu / danych & Okres przechowywania & Podstawa prawna \\
\midrule\noalign{}
\endhead
\bottomrule\noalign{}
\endlastfoot
Księgi rachunkowe i dokumenty finansowe & 5 lat od zakończenia roku
obrotowego & Ustawa o rachunkowości art. 74 \\
Deklaracje podatkowe & 5 lat od upływu terminu płatności podatku &
Ordynacja podatkowa art. 70 \\
Umowy z klientami & 10 lat od wygaśnięcia umowy & Kodeks cywilny --
przedawnienie \\
Akta pracownicze & 10 lat (pracownicy od 01.01.2019) & Ustawa
archiwalna \\
Dane marketingowe (zgoda) & Do cofnięcia zgody + 1 rok & RODO art. 7 \\
Logi systemowe i dostępu & 12 miesięcy & Polityka wewnętrzna \\
Kopie zapasowe (backup) & 90 dni (rotacyjnie) & Polityka wewnętrzna \\
\end{longtable}

\subsection{4. Prawa osób, których dane
dotyczą}\label{prawa-osuxf3b-ktuxf3rych-dane-dotyczux105}

\begin{itemize}
\tightlist
\item
  \textbf{Prawo dostępu} (art. 15 RODO) -- odpowiedź w ciągu 30 dni
\item
  \textbf{Prawo do sprostowania} (art. 16 RODO)
\item
  \textbf{Prawo do usunięcia} „prawo do bycia zapomnianym" (art. 17
  RODO)
\item
  \textbf{Prawo do ograniczenia przetwarzania} (art. 18 RODO)
\item
  \textbf{Prawo do przenoszenia danych} (art. 20 RODO)
\item
  \textbf{Prawo do sprzeciwu} (art. 21 RODO)
\end{itemize}

Punkt kontaktowy dla wniosków: \textbf{{[}EMAIL/ADRES{]}} \textbar{}
Termin rozpatrzenia: 30 dni (możliwość przedłużenia o 60 dni).

\subsection{5. Bezpieczeństwo danych}\label{bezpieczeux144stwo-danych}

\begin{itemize}
\tightlist
\item
  Dane przetwarzane wyłącznie na zatwierdzonych urządzeniach i
  platformach
\item
  Szyfrowanie danych w spoczynku (AES-256) i podczas transmisji (TLS
  1.2+)
\item
  Dostęp na zasadzie „najmniejszych uprawnień" (least privilege)
\item
  Regularne kopie zapasowe wg polityki backup
\item
  Pseudonimizacja danych tam, gdzie to możliwe
\end{itemize}

\subsection{6. Naruszenie ochrony danych --
procedura}\label{naruszenie-ochrony-danych-procedura}

\begin{enumerate}
\tightlist
\item
  Wykrycie naruszenia → natychmiastowe poinformowanie przełożonego /
  właściciela
\item
  Ocena ryzyka -- czy naruszenie stwarza ryzyko dla praw osób?
\item
  \textbf{Zgłoszenie do UODO} -- w ciągu \textbf{72 godzin} od wykrycia
  (jeśli ryzyko wystąpiło)
\item
  Powiadomienie osób poszkodowanych -- bez zbędnej zwłoki (jeśli wysokie
  ryzyko)
\item
  Dokumentacja incydentu -- rejestr naruszeń (wewnętrzny)
\end{enumerate}

\subsection{7. Powierzenie danych podmiotom
zewnętrznym}\label{powierzenie-danych-podmiotom-zewnux119trznym}

Każde powierzenie danych wymaga zawarcia \textbf{umowy powierzenia
przetwarzania (DPA)}. Lista zatwierdzonych procesorów:

\begin{longtable}[]{@{}>{\raggedright\arraybackslash}p{\dimexpr\linewidth/1000*270-2\tabcolsep\relax}>{\raggedright\arraybackslash}p{\dimexpr\linewidth/1000*430-2\tabcolsep\relax}>{\raggedright\arraybackslash}p{\dimexpr\linewidth/1000*300-2\tabcolsep\relax}@{}}
\toprule\noalign{}
Procesor & Zakres przetwarzania & Umowa DPA \\
\midrule\noalign{}
\endhead
\bottomrule\noalign{}
\endlastfoot
{[}NAZWA DOSTAWCY CHMURY{]} & Przechowywanie plików & {[} {]} TAK ~{[}
{]} NIE ~Data: \underline{\hspace{2cm}} \\
{[}BIURO IT / ZEWNĘTRZNY IT{]} & Wsparcie informatyczne & {[} {]} TAK
~{[} {]} NIE ~Data: \underline{\hspace{2cm}} \\
{[}INNE{]} & \underline{\hspace{3cm}} & {[} {]} TAK ~{[}
{]} NIE ~Data: \underline{\hspace{2cm}} \\
\end{longtable}

\begin{longtable}[]{@{}
  >{\raggedright\arraybackslash}p{\dimexpr\linewidth/1000*500-1\tabcolsep\relax}
  >{\raggedright\arraybackslash}p{\dimexpr\linewidth/1000*500-1\tabcolsep\relax}@{}}
\toprule\noalign{}
\endhead
\bottomrule\noalign{}
\endlastfoot
\begin{minipage}[t]{\linewidth}\raggedright
Zatwierdził:

Podpis: \underline{\hspace{3cm}}

Data: \underline{\hspace{4cm}}
\end{minipage} & \begin{minipage}[t]{\linewidth}\raggedright
Stanowisko / Tytuł:

Imię i nazwisko: \underline{\hspace{3cm}}

Funkcja: \underline{\hspace{3cm}}
\end{minipage} \\
\end{longtable}

\phantomsection\label{doc2}
PROC-002

\section{PROCEDURA WERYFIKACJI CERTYFIKACJI I BEZPIECZEŃSTWA
SPRZĘTU}\label{procedura-weryfikacji-certyfikacji-i-bezpieczeux144stwa-sprzux119tu}

\begin{longtable}[]{@{}>{\raggedright\arraybackslash}p{\dimexpr\linewidth/1000*300-2\tabcolsep\relax}>{\raggedright\arraybackslash}p{\dimexpr\linewidth/1000*700-2\tabcolsep\relax}@{}}
\toprule\noalign{}
\endhead
\bottomrule\noalign{}
\endlastfoot
Wersja: & 1.0 \\
Dotyczy: & Komputery, laptopy, switche, routery, drukarki, UPS \\
Odpowiedzialny: & Administrator IT / Właściciel \\
Częstotliwość: & Przy zakupie + coroczna weryfikacja \\
\end{longtable}

\subsection{1. Wymagania przy zakupie nowego
sprzętu}\label{wymagania-przy-zakupie-nowego-sprzux119tu}

\subsubsection{1.1 Komputery i laptopy}\label{komputery-i-laptopy}

\begin{itemize}
\tightlist
\item
  Zakup wyłącznie od autoryzowanych dystrybutorów (faktura VAT,
  gwarancja producenta)
\item
  Wymagane certyfikaty: CE, FCC, RoHS
\item
  Wymagania minimalne: TPM 2.0, Secure Boot, szyfrowanie dysku
  (BitLocker/FileVault)
\item
  Zatwierdzenie przez administratora IT przed zakupem
\item
  Rejestracja w inwentarzu sprzętowym -- numer seryjny, data zakupu,
  użytkownik
\end{itemize}

\subsubsection{1.2 Routery i switche}\label{routery-i-switche}

\begin{itemize}
\tightlist
\item
  Zakup od renomowanych producentów (Cisco, Fortinet, Ubiquiti, TP-Link
  Business)
\item
  Zmiana domyślnych haseł administratora -- bezpośrednio po rozpakowaniu
\item
  Wyłączenie zdalnego zarządzania przez WAN (jeśli niepotrzebne)
\item
  Włączenie automatycznych aktualizacji firmware lub harmonogram
  ręcznych aktualizacji
\item
  Dokumentacja konfiguracji sieciowej przechowywana offline
  (zaszyfrowana)
\end{itemize}

\subsubsection{1.3 Drukarki i urządzenia
wielofunkcyjne}\label{drukarki-i-urzux105dzenia-wielofunkcyjne}

\begin{itemize}
\tightlist
\item
  Zmiana domyślnych danych logowania administratora
\item
  Wyłączenie usług sieciowych nieużywanych (FTP, Telnet, HTTP → wymóg
  HTTPS)
\item
  Szyfrowanie połączenia (TLS) dla drukowania sieciowego
\item
  Włączenie bezpiecznego wydruku (PIN przed wydrukiem)
\item
  Regularne czyszczenie pamięci wewnętrznej drukarki
\end{itemize}

\subsection{2. Harmonogram weryfikacji -- lista
kontrolna}\label{harmonogram-weryfikacji-lista-kontrolna}

\begin{longtable}[]{@{}>{\raggedright\arraybackslash}p{\dimexpr\linewidth/1000*270-2\tabcolsep\relax}>{\raggedright\arraybackslash}p{\dimexpr\linewidth/1000*430-2\tabcolsep\relax}>{\raggedright\arraybackslash}p{\dimexpr\linewidth/1000*300-2\tabcolsep\relax}@{}}
\toprule\noalign{}
Czynność weryfikacyjna & Częstotliwość & Status {[} {]} / Data \\
\midrule\noalign{}
\endhead
\bottomrule\noalign{}
\endlastfoot
Aktualizacja firmware routerów i switchy & Co 3 miesiące & {[} {]}
\underline{\hspace{2cm}} \\
Weryfikacja aktywnego oprogramowania antywirusowego & Co miesiąc & {[}
{]} \underline{\hspace{2cm}} \\
Skan podatności sieci & Co 6 miesięcy & {[} {]}
\underline{\hspace{2cm}} \\
Przegląd listy urządzeń podłączonych do sieci & Co miesiąc & {[} {]}
\underline{\hspace{2cm}} \\
Weryfikacja szyfrowania dysków (BitLocker/FileVault) & Co 6 miesięcy &
{[} {]} \underline{\hspace{2cm}} \\
Aktualizacja inwentarza sprzętowego & Co 12 miesięcy & {[} {]}
\underline{\hspace{2cm}} \\
Przegląd konfiguracji firewall & Co 6 miesięcy & {[} {]}
\underline{\hspace{2cm}} \\
\end{longtable}

\subsection{3. Rejestr sprzętu
(szablon)}\label{rejestr-sprzux119tu-szablon}

\begin{longtable}[]{@{}>{\raggedright\arraybackslash}p{\dimexpr\linewidth/1000*180-2\tabcolsep\relax}>{\raggedright\arraybackslash}p{\dimexpr\linewidth/1000*170-2\tabcolsep\relax}>{\raggedright\arraybackslash}p{\dimexpr\linewidth/1000*170-2\tabcolsep\relax}>{\raggedright\arraybackslash}p{\dimexpr\linewidth/1000*170-2\tabcolsep\relax}>{\raggedright\arraybackslash}p{\dimexpr\linewidth/1000*170-2\tabcolsep\relax}>{\raggedright\arraybackslash}p{\dimexpr\linewidth/1000*140-2\tabcolsep\relax}@{}}
\toprule\noalign{}
Urządzenie & Nr seryjny & Użytkownik & Data zakupu & Ostatni przegląd &
Status \\
\midrule\noalign{}
\endhead
\bottomrule\noalign{}
\endlastfoot
\underline{\hspace{2cm}} & \underline{\hspace{2cm}} &
\underline{\hspace{2cm}} & \underline{\hspace{2cm}} &
\underline{\hspace{2cm}} & {[} {]} OK \\
\underline{\hspace{2cm}} & \underline{\hspace{2cm}} &
\underline{\hspace{2cm}} & \underline{\hspace{2cm}} &
\underline{\hspace{2cm}} & {[} {]} OK \\
\underline{\hspace{2cm}} & \underline{\hspace{2cm}} &
\underline{\hspace{2cm}} & \underline{\hspace{2cm}} &
\underline{\hspace{2cm}} & {[} {]} OK \\
\end{longtable}

\begin{longtable}[]{@{}
  >{\raggedright\arraybackslash}p{\dimexpr\linewidth/1000*500-1\tabcolsep\relax}
  >{\raggedright\arraybackslash}p{\dimexpr\linewidth/1000*500-1\tabcolsep\relax}@{}}
\toprule\noalign{}
\endhead
\bottomrule\noalign{}
\endlastfoot
\begin{minipage}[t]{\linewidth}\raggedright
Zatwierdził:

Podpis: \underline{\hspace{3cm}}

Data: \underline{\hspace{4cm}}
\end{minipage} & \begin{minipage}[t]{\linewidth}\raggedright
Stanowisko / Tytuł:

Imię i nazwisko: \underline{\hspace{3cm}}

Funkcja: \underline{\hspace{3cm}}
\end{minipage} \\
\end{longtable}

\phantomsection\label{doc3}
PROC-003

\section{PROCEDURA WERYFIKACJI CERTYFIKACJI I BEZPIECZEŃSTWA
OPROGRAMOWANIA}\label{procedura-weryfikacji-certyfikacji-i-bezpieczeux144stwa-oprogramowania}

\begin{longtable}[]{@{}>{\raggedright\arraybackslash}p{\dimexpr\linewidth/1000*300-2\tabcolsep\relax}>{\raggedright\arraybackslash}p{\dimexpr\linewidth/1000*700-2\tabcolsep\relax}@{}}
\toprule\noalign{}
\endhead
\bottomrule\noalign{}
\endlastfoot
Wersja: & 1.0 \\
Dotyczy: & Systemy operacyjne, aplikacje biznesowe, oprogramowanie
księgowe, przeglądarki \\
Odpowiedzialny: & Administrator IT \\
Częstotliwość: & Przy instalacji + kwartalny przegląd \\
\end{longtable}

\subsection{1. Lista Zatwierdzonego Oprogramowania
(LZO)}\label{lista-zatwierdzonego-oprogramowania-lzo}

W firmie dopuszczone jest wyłącznie oprogramowanie z LZO. Instalacja
oprogramowania spoza listy wymaga pisemnej zgody administratora IT.

\begin{longtable}[]{@{}>{\raggedright\arraybackslash}p{\dimexpr\linewidth/1000*270-2\tabcolsep\relax}>{\raggedright\arraybackslash}p{\dimexpr\linewidth/1000*250-2\tabcolsep\relax}>{\raggedright\arraybackslash}p{\dimexpr\linewidth/1000*250-2\tabcolsep\relax}>{\raggedright\arraybackslash}p{\dimexpr\linewidth/1000*230-2\tabcolsep\relax}@{}}
\toprule\noalign{}
Kategoria & Zatwierdzone rozwiązanie & Cel & Licencja OK \\
\midrule\noalign{}
\endhead
\bottomrule\noalign{}
\endlastfoot
System operacyjny & Windows 11 Pro / macOS 13+ & Stacje robocze & {[}
{]} TAK \\
Oprogramowanie księgowe & {[}NAZWA -- np. Insert/Symfonia{]} & Główna
aplikacja & {[} {]} TAK \\
Antywirus/EDR & {[}NAZWA{]} & Ochrona endpoint & {[} {]} TAK \\
VPN & {[}NAZWA{]} & Zdalny dostęp & {[} {]} TAK \\
Przeglądarka & Chrome / Firefox (aktualna wersja) & Internet & {[} {]}
TAK \\
Office & Microsoft 365 Business & Praca biurowa & {[} {]} TAK \\
Manager haseł & {[}NAZWA -- np. Bitwarden/1Password{]} & Zarządzanie
hasłami & {[} {]} TAK \\
\end{longtable}

\subsection{2. Procedura instalacji nowego
oprogramowania}\label{procedura-instalacji-nowego-oprogramowania}

\begin{enumerate}
\tightlist
\item
  Wniosek pracownika do administratora IT (forma pisemna/email)
\item
  Ocena bezpieczeństwa: źródło, producent, recenzje, sprawdzenie CVE
\item
  Weryfikacja licencji -- oprogramowanie komercyjne lub Open Source z
  zatwierdzoną licencją
\item
  Instalacja przez administratora lub z jego zgodą (nie przez
  użytkownika)
\item
  Dodanie do LZO i dokumentacja
\item
  Monitorowanie przez 30 dni po instalacji
\end{enumerate}

\subsection{3. Aktualizacje i łatki
bezpieczeństwa}\label{aktualizacje-i-ux142atki-bezpieczeux144stwa}

\begin{itemize}
\tightlist
\item
  Aktualizacje systemu operacyjnego: automatyczne lub w ciągu \textbf{72
  godzin} od wydania łatki krytycznej
\item
  Oprogramowanie biznesowe: aktualizacja w ciągu \textbf{14 dni} od
  wydania
\item
  Przeglądarki i wtyczki: automatyczne aktualizacje włączone
\item
  Oprogramowanie bez aktualizacji od \textgreater12 miesięcy -- ocena i
  ewentualne odinstalowanie
\end{itemize}

\subsection{4. Lista kontrolna -- kwartalny przegląd
oprogramowania}\label{lista-kontrolna-kwartalny-przeglux105d-oprogramowania}

\begin{longtable}[]{@{}>{\raggedright\arraybackslash}p{\dimexpr\linewidth/1000*270-2\tabcolsep\relax}>{\raggedright\arraybackslash}p{\dimexpr\linewidth/1000*430-2\tabcolsep\relax}>{\raggedright\arraybackslash}p{\dimexpr\linewidth/1000*300-2\tabcolsep\relax}@{}}
\toprule\noalign{}
Czynność & Termin & Status \\
\midrule\noalign{}
\endhead
\bottomrule\noalign{}
\endlastfoot
Weryfikacja wersji systemu operacyjnego na wszystkich komputerach & Co
kwartał & {[} {]} \underline{\hspace{2cm}} \\
Przegląd listy zainstalowanego oprogramowania (porównanie z LZO) & Co
kwartał & {[} {]} \underline{\hspace{2cm}} \\
Sprawdzenie dat wygaśnięcia licencji & Co kwartał & {[} {]}
\underline{\hspace{2cm}} \\
Weryfikacja bazy sygnatur antywirusa & Co miesiąc & {[} {]}
\underline{\hspace{2cm}} \\
Skan podatności zainstalowanego oprogramowania (CVE) & Co 6 miesięcy &
{[} {]} \underline{\hspace{2cm}} \\
Usunięcie nieautoryzowanego lub nieaktualnego SW & Bieżąco & {[} {]}
\underline{\hspace{2cm}} \\
\end{longtable}

\begin{longtable}[]{@{}
  >{\raggedright\arraybackslash}p{\dimexpr\linewidth/1000*500-1\tabcolsep\relax}
  >{\raggedright\arraybackslash}p{\dimexpr\linewidth/1000*500-1\tabcolsep\relax}@{}}
\toprule\noalign{}
\endhead
\bottomrule\noalign{}
\endlastfoot
\begin{minipage}[t]{\linewidth}\raggedright
Zatwierdził:

Podpis: \underline{\hspace{3cm}}

Data: \underline{\hspace{4cm}}
\end{minipage} & \begin{minipage}[t]{\linewidth}\raggedright
Stanowisko / Tytuł:

Imię i nazwisko: \underline{\hspace{3cm}}

Funkcja: \underline{\hspace{3cm}}
\end{minipage} \\
\end{longtable}

\phantomsection\label{doc4}
PROC-004

\section{PROCEDURA ZABEZPIECZANIA INFRASTRUKTURY PRZED
ATAKAMI}\label{procedura-zabezpieczania-infrastruktury-przed-atakami}

\begin{longtable}[]{@{}>{\raggedright\arraybackslash}p{\dimexpr\linewidth/1000*300-2\tabcolsep\relax}>{\raggedright\arraybackslash}p{\dimexpr\linewidth/1000*700-2\tabcolsep\relax}@{}}
\toprule\noalign{}
\endhead
\bottomrule\noalign{}
\endlastfoot
Wersja: & 1.0 \\
Dotyczy: & Sieć firmowa, serwery, stacje robocze, dostęp zdalny \\
Odpowiedzialny: & Administrator IT \\
Przegląd: & Co 6 miesięcy \\
\end{longtable}

\subsection{1. Architektura sieciowa}\label{architektura-sieciowa}

\begin{itemize}
\tightlist
\item
  Segmentacja sieci: sieć biurowa (LAN), sieć dla gości (Guest WiFi),
  VLAN dla serwerów -- odseparowane
\item
  Firewall na granicy sieci -- polityka „odmów wszystkiego, zezwalaj na
  wyjątki"
\item
  Brak bezpośredniego dostępu RDP/SSH z Internetu bez VPN
\item
  Dokumentacja topologii sieci aktualizowana przy każdej zmianie
\end{itemize}

\subsection{2. Hasła i
uwierzytelnianie}\label{hasux142a-i-uwierzytelnianie}

\begin{itemize}
\tightlist
\item
  Minimalna długość hasła: \textbf{12 znaków}
\item
  Złożoność: wielkie/małe litery, cyfry, znaki specjalne
\item
  Zmiana haseł: co 90 dni (konta uprzywilejowane) / 180 dni
  (standardowe)
\item
  Obowiązkowe \textbf{MFA} dla: kont administratora, VPN, poczty e-mail,
  systemów księgowych
\item
  Zarządzanie hasłami przez zatwierdzonego managera haseł (nie Excel,
  nie kartka!)
\item
  Blokada konta po \textbf{5 błędnych próbach} logowania
\end{itemize}

\subsection{3. Dostęp zdalny}\label{dostux119p-zdalny}

\begin{itemize}
\tightlist
\item
  Zdalny dostęp wyłącznie przez VPN (z MFA)
\item
  Lista zatwierdzonych użytkowników VPN -- przegląd co 3 miesiące
\item
  Sesje VPN z limitem czasu (timeout po 30 min bezczynności)
\item
  Zakaz używania publicznych sieci WiFi do pracy z danymi klientów bez
  VPN
\end{itemize}

\subsection{4. Ochrona przed złośliwym
oprogramowaniem}\label{ochrona-przed-zux142oux15bliwym-oprogramowaniem}

\begin{itemize}
\tightlist
\item
  Oprogramowanie EDR/AV na każdym urządzeniu końcowym -- obowiązkowo
\item
  Blokowanie wykonywania skryptów z makrami w dokumentach Office
  (polityka grupowa)
\item
  Filtrowanie DNS z blokowaniem złośliwych domen
\item
  Polityka białej listy aplikacji (Application Whitelisting) na stacjach
  roboczych
\item
  Skanowanie załączników e-mail przez antywirus
\end{itemize}

\subsection{5. Polityka kopii zapasowych
(3-2-1)}\label{polityka-kopii-zapasowych-3-2-1}

Reguła 3-2-1: \textbf{3 kopie}, na \textbf{2 różnych nośnikach}, w tym
\textbf{1 offsite}. Testowanie przywracania: co kwartał.

\begin{longtable}[]{@{}>{\raggedright\arraybackslash}p{\dimexpr\linewidth/1000*270-2\tabcolsep\relax}>{\raggedright\arraybackslash}p{\dimexpr\linewidth/1000*250-2\tabcolsep\relax}>{\raggedright\arraybackslash}p{\dimexpr\linewidth/1000*250-2\tabcolsep\relax}>{\raggedright\arraybackslash}p{\dimexpr\linewidth/1000*230-2\tabcolsep\relax}@{}}
\toprule\noalign{}
Rodzaj backup & Częstotliwość & Przechowywanie & Szyfrowanie \\
\midrule\noalign{}
\endhead
\bottomrule\noalign{}
\endlastfoot
Pełny backup serwerów & Co tydzień (niedziela) & 90 dni & {[} {]} TAK \\
Przyrostowy backup plików & Codziennie & 30 dni & {[} {]} TAK \\
Backup baz danych (księgowość) & Codziennie (noc) & 90 dni & {[} {]}
TAK \\
Backup poza siedzibą (offsite) & Co tydzień & 12 miesięcy & {[} {]}
TAK \\
\end{longtable}

\begin{longtable}[]{@{}
  >{\raggedright\arraybackslash}p{\dimexpr\linewidth/1000*500-1\tabcolsep\relax}
  >{\raggedright\arraybackslash}p{\dimexpr\linewidth/1000*500-1\tabcolsep\relax}@{}}
\toprule\noalign{}
\endhead
\bottomrule\noalign{}
\endlastfoot
\begin{minipage}[t]{\linewidth}\raggedright
Zatwierdził:

Podpis: \underline{\hspace{3cm}}

Data: \underline{\hspace{4cm}}
\end{minipage} & \begin{minipage}[t]{\linewidth}\raggedright
Stanowisko / Tytuł:

Imię i nazwisko: \underline{\hspace{3cm}}

Funkcja: \underline{\hspace{3cm}}
\end{minipage} \\
\end{longtable}

\phantomsection\label{doc5}
POL-005

\section{POLITYKA DOSTĘPU I MONITOROWANIA
UŻYTKOWNIKÓW}\label{polityka-dostux119pu-i-monitorowania-uux17cytkownikuxf3w}

\begin{longtable}[]{@{}>{\raggedright\arraybackslash}p{\dimexpr\linewidth/1000*300-2\tabcolsep\relax}>{\raggedright\arraybackslash}p{\dimexpr\linewidth/1000*700-2\tabcolsep\relax}@{}}
\toprule\noalign{}
\endhead
\bottomrule\noalign{}
\endlastfoot
Wersja: & 1.0 \\
Dotyczy: & Wszyscy użytkownicy systemów informatycznych firmy \\
Odpowiedzialny: & Administrator IT / Kierownik \\
Przegląd: & Co 12 miesięcy lub przy zmianie struktury organizacyjnej \\
\end{longtable}

\subsection{1. Zasada najniższych
uprawnień}\label{zasada-najniux17cszych-uprawnieux144}

Każdy użytkownik otrzymuje dostęp wyłącznie do zasobów niezbędnych do
wykonywania swoich obowiązków. Uprawnienia są nadawane przez
administratora na pisemny wniosek kierownika.

\subsection{2. Macierz dostępu --
szablon}\label{macierz-dostux119pu-szablon}

\begin{longtable}[]{@{}>{\raggedright\arraybackslash}p{\dimexpr\linewidth/1000*220-2\tabcolsep\relax}>{\raggedright\arraybackslash}p{\dimexpr\linewidth/1000*200-2\tabcolsep\relax}>{\raggedright\arraybackslash}p{\dimexpr\linewidth/1000*200-2\tabcolsep\relax}>{\raggedright\arraybackslash}p{\dimexpr\linewidth/1000*200-2\tabcolsep\relax}>{\raggedright\arraybackslash}p{\dimexpr\linewidth/1000*180-2\tabcolsep\relax}@{}}
\toprule\noalign{}
Zasób / System & Właściciel & Księgowy & Asystent & Gość/Klient \\
\midrule\noalign{}
\endhead
\bottomrule\noalign{}
\endlastfoot
System księgowy (pełny) & Pełny & Pełny & Odczyt & Brak \\
Dane klientów (CRM) & Pełny & Odczyt/Zapis & Odczyt & Brak \\
Dokumenty finansowe (serwer) & Pełny & Swoje teczki & Brak & Brak \\
Poczta e-mail firmowa & Pełny & Własna skrzynka & Własna skrzynka &
Brak \\
Internet / VPN & Pełny & Tak & Ograniczony & Brak \\
Panel administracyjny IT & Pełny & Brak & Brak & Brak \\
\end{longtable}

\subsection{3. Monitorowanie działań
użytkowników}\label{monitorowanie-dziaux142aux144-uux17cytkownikuxf3w}

Firma zastrzega sobie prawo do monitorowania korzystania z systemów
informatycznych. Użytkownicy są poinformowani o fakcie monitorowania
(klauzula w umowie i regulaminie IT).

Monitorowane działania obejmują: logowania i wylogowania, dostęp do
plików i baz danych, wysyłane e-maile z firmowych kont, dostęp do
Internetu, drukowanie dokumentów, podpinane urządzenia zewnętrzne (USB).

Dane z monitorowania są przechowywane przez 12 miesięcy i dostępne
wyłącznie dla administratora IT i właściciela firmy.

\subsection{4. Akceptowalne użycie zasobów
IT}\label{akceptowalne-uux17cycie-zasobuxf3w-it}

\subsubsection{Dozwolone:}\label{dozwolone}

\begin{itemize}
\tightlist
\item
  Używanie sprzętu i oprogramowania do celów służbowych
\item
  Ograniczone użycie prywatne niezakłócające pracy
\end{itemize}

\subsubsection{Zabronione:}\label{zabronione}

\begin{itemize}
\tightlist
\item
  Instalacja nieautoryzowanego oprogramowania
\item
  Podłączanie prywatnych nośników USB bez zgody administratora
\item
  Udostępnianie danych logowania innym osobom
\item
  Obchodzenie zabezpieczeń (wyłączanie antywirusa, omijanie VPN)
\item
  Przesyłanie danych klientów na prywatne skrzynki e-mail / chmury
\item
  Używanie służbowych kont do celów prywatnych / sprzedażowych
\end{itemize}

\begin{longtable}[]{@{}
  >{\raggedright\arraybackslash}p{\dimexpr\linewidth/1000*500-1\tabcolsep\relax}
  >{\raggedright\arraybackslash}p{\dimexpr\linewidth/1000*500-1\tabcolsep\relax}@{}}
\toprule\noalign{}
\endhead
\bottomrule\noalign{}
\endlastfoot
\begin{minipage}[t]{\linewidth}\raggedright
Zatwierdził:

Podpis: \underline{\hspace{3cm}}

Data: \underline{\hspace{4cm}}
\end{minipage} & \begin{minipage}[t]{\linewidth}\raggedright
Stanowisko / Tytuł:

Imię i nazwisko: \underline{\hspace{3cm}}

Funkcja: \underline{\hspace{3cm}}
\end{minipage} \\
\end{longtable}

\phantomsection\label{doc6}
PROC-006

\section{PROCEDURA ONBOARDINGU NOWYCH UŻYTKOWNIKÓW
(IT)}\label{procedura-onboardingu-nowych-uux17cytkownikuxf3w-it}

\begin{longtable}[]{@{}>{\raggedright\arraybackslash}p{\dimexpr\linewidth/1000*300-2\tabcolsep\relax}>{\raggedright\arraybackslash}p{\dimexpr\linewidth/1000*700-2\tabcolsep\relax}@{}}
\toprule\noalign{}
\endhead
\bottomrule\noalign{}
\endlastfoot
Wersja: & 1.0 \\
Dotyczy: & Nowi pracownicy, stażyści, współpracownicy \\
Odpowiedzialny: & Administrator IT + HR \\
Termin realizacji: & Dzień przed lub w dniu pierwszego dnia pracy \\
\end{longtable}

\subsection{1. Lista kontrolna onboardingu
IT}\label{lista-kontrolna-onboardingu-it}

\begin{longtable}[]{@{}>{\raggedright\arraybackslash}p{\dimexpr\linewidth/1000*270-2\tabcolsep\relax}>{\raggedright\arraybackslash}p{\dimexpr\linewidth/1000*250-2\tabcolsep\relax}>{\raggedright\arraybackslash}p{\dimexpr\linewidth/1000*250-2\tabcolsep\relax}>{\raggedright\arraybackslash}p{\dimexpr\linewidth/1000*230-2\tabcolsep\relax}@{}}
\toprule\noalign{}
\# & Czynność & Odpowiedzialny & Status / Data \\
\midrule\noalign{}
\endhead
\bottomrule\noalign{}
\endlastfoot
1 & Odebranie wypełnionego wniosku o konto od kierownika & HR /
Kierownik & {[} {]} \underline{\hspace{2cm}} \\
2 & Utworzenie konta użytkownika w AD / systemie & Administrator IT &
{[} {]} \underline{\hspace{2cm}} \\
3 & Nadanie uprawnień wg macierzy dostępu (POL-005) & Administrator IT &
{[} {]} \underline{\hspace{2cm}} \\
4 & Konfiguracja MFA (telefon / aplikacja) & Administrator IT & {[} {]}
\underline{\hspace{2cm}} \\
5 & Wydanie sprzętu -- laptop, akcesoria (protokół zdawczo-odbiorczy) &
Administrator IT & {[} {]} \underline{\hspace{2cm}} \\
6 & Instalacja i konfiguracja wymaganego oprogramowania & Administrator
IT & {[} {]} \underline{\hspace{2cm}} \\
7 & Konfiguracja poczty e-mail i podpisu & Administrator IT & {[} {]}
\underline{\hspace{2cm}} \\
8 & Podłączenie do VPN + test połączenia & Administrator IT & {[} {]}
\underline{\hspace{2cm}} \\
9 & Zapoznanie z Polityką Bezpieczeństwa IT (POL-001, POL-005) &
Administrator IT & {[} {]} \underline{\hspace{2cm}} \\
10 & Szkolenie wprowadzające z cyberbezpieczeństwa (PROC-008) &
Administrator IT & {[} {]} \underline{\hspace{2cm}} \\
11 & Podpisanie Oświadczenia o zapoznaniu z politykami IT & HR /
Pracownik & {[} {]} \underline{\hspace{2cm}} \\
12 & Podpisanie Oświadczenia o poufności (NDA / klauzula RODO) & HR /
Pracownik & {[} {]} \underline{\hspace{2cm}} \\
13 & Test dostępu -- potwierdzenie że wszystko działa & Nowy pracownik &
{[} {]} \underline{\hspace{2cm}} \\
\end{longtable}

\subsection{2. Oświadczenie o zapoznaniu z politykami IT (do
podpisania)}\label{oux15bwiadczenie-o-zapoznaniu-z-politykami-it-do-podpisania}

Ja, niżej podpisany/a {~}, potwierdzam, że:

\begin{enumerate}
\tightlist
\item
  Zapoznałem/am się z Polityką Bezpieczeństwa IT firmy {[}NAZWA{]}
\item
  Zostałem/am poinformowany/a o fakcie monitorowania aktywności w
  systemach IT
\item
  Zobowiązuję się do przestrzegania zasad określonych w politykach
\item
  Przyjmuję do wiadomości, że naruszenie polityk może skutkować
  konsekwencjami służbowymi
\end{enumerate}

Data: \underline{\hspace{3cm}}

Podpis pracownika: \underline{\hspace{3cm}}

Podpis przełożonego: \underline{\hspace{3cm}}

\begin{longtable}[]{@{}
  >{\raggedright\arraybackslash}p{\dimexpr\linewidth/1000*500-1\tabcolsep\relax}
  >{\raggedright\arraybackslash}p{\dimexpr\linewidth/1000*500-1\tabcolsep\relax}@{}}
\toprule\noalign{}
\endhead
\bottomrule\noalign{}
\endlastfoot
\begin{minipage}[t]{\linewidth}\raggedright
Zatwierdził:

Podpis: \underline{\hspace{3cm}}

Data: \underline{\hspace{4cm}}
\end{minipage} & \begin{minipage}[t]{\linewidth}\raggedright
Stanowisko / Tytuł:

Imię i nazwisko: \underline{\hspace{3cm}}

Funkcja: \underline{\hspace{3cm}}
\end{minipage} \\
\end{longtable}

\phantomsection\label{doc7}
PROC-007

\section{PROCEDURA OFFBOARDINGU ODCHODZĄCYCH UŻYTKOWNIKÓW
(IT)}\label{procedura-offboardingu-odchodzux105cych-uux17cytkownikuxf3w-it}

\begin{longtable}[]{@{}>{\raggedright\arraybackslash}p{\dimexpr\linewidth/1000*300-2\tabcolsep\relax}>{\raggedright\arraybackslash}p{\dimexpr\linewidth/1000*700-2\tabcolsep\relax}@{}}
\toprule\noalign{}
\endhead
\bottomrule\noalign{}
\endlastfoot
Wersja: & 1.0 \\
Dotyczy: & Pracownicy, stażyści, współpracownicy kończący współpracę \\
Odpowiedzialny: & Administrator IT + HR \\
Termin realizacji: & W dniu lub natychmiast po ostatnim dniu pracy \\
\end{longtable}

\textbf{UWAGA:} Konta należy dezaktywować w dniu zakończenia współpracy -- nie
wcześniej, nie później (ryzyko wycieku po odejściu lub brak dostępu w
ostatnim dniu).

\subsection{1. Lista kontrolna offboardingu
IT}\label{lista-kontrolna-offboardingu-it}

\begin{longtable}[]{@{}>{\raggedright\arraybackslash}p{\dimexpr\linewidth/1000*270-2\tabcolsep\relax}>{\raggedright\arraybackslash}p{\dimexpr\linewidth/1000*250-2\tabcolsep\relax}>{\raggedright\arraybackslash}p{\dimexpr\linewidth/1000*250-2\tabcolsep\relax}>{\raggedright\arraybackslash}p{\dimexpr\linewidth/1000*230-2\tabcolsep\relax}@{}}
\toprule\noalign{}
\# & Czynność & Odpowiedzialny & Status / Data \\
\midrule\noalign{}
\endhead
\bottomrule\noalign{}
\endlastfoot
1 & Wyłączenie konta użytkownika (nie usuwanie -- archiwum 90 dni) &
Administrator IT & {[} {]} \underline{\hspace{2cm}} \\
2 & Zmiana haseł do współdzielonych kont (jeśli dotyczy) & Administrator
IT & {[} {]} \underline{\hspace{2cm}} \\
3 & Odwołanie dostępu VPN & Administrator IT & {[} {]}
\underline{\hspace{2cm}} \\
4 & Usunięcie / dezaktywacja tokena MFA & Administrator IT & {[} {]}
\underline{\hspace{2cm}} \\
5 & Archiwizacja skrzynki e-mail (90 dni przekierowania do przełożonego)
& Administrator IT & {[} {]} \underline{\hspace{2cm}} \\
6 & Przejęcie plików służbowych (przekazanie kierownikowi) & Kierownik &
{[} {]} \underline{\hspace{2cm}} \\
7 & Zwrot sprzętu (laptop, telefon, klucze, karty dostępu) -- protokół &
HR & {[} {]} \underline{\hspace{2cm}} \\
8 & Wyczyszczenie danych z urządzeń zwracanych (certyfikowane usunięcie)
& Administrator IT & {[} {]} \underline{\hspace{2cm}} \\
9 & Usunięcie z list mailingowych i grup & Administrator IT & {[} {]}
\underline{\hspace{2cm}} \\
10 & Odwołanie dostępu do usług zewnętrznych (chmura, narzędzia SaaS) &
Administrator IT & {[} {]} \underline{\hspace{2cm}} \\
11 & Podpisanie oświadczenia o zwrocie danych i poufności & HR /
Pracownik & {[} {]} \underline{\hspace{2cm}} \\
12 & Aktualizacja macierzy dostępu i inwentarza sprzętowego &
Administrator IT & {[} {]} \underline{\hspace{2cm}} \\
\end{longtable}

\subsection{2. Oświadczenie o zwrocie danych i poufności
(Offboarding)}\label{oux15bwiadczenie-o-zwrocie-danych-i-poufnoux15bci-offboarding}

Ja, niżej podpisany/a {~}, potwierdzam, że:

\begin{enumerate}
\tightlist
\item
  Zwróciłem/am cały sprzęt firmowy (wymieniony w załączniku protokołu)
\item
  Usunąłem/am wszelkie dane firmowe z urządzeń prywatnych
\item
  Nie posiadam kopii danych klientów, dokumentów finansowych ani haseł
  firmowych
\item
  Zobowiązuję się do zachowania poufności przez {[}3/5{]} lat po
  zakończeniu współpracy
\item
  Jestem świadomy/a konsekwencji prawnych naruszenia powyższych
  zobowiązań
\end{enumerate}

Data: \underline{\hspace{3cm}}

Podpis pracownika: \underline{\hspace{3cm}}

Podpis HR / przełożonego: \underline{\hspace{3cm}}

\begin{longtable}[]{@{}
  >{\raggedright\arraybackslash}p{\dimexpr\linewidth/1000*500-1\tabcolsep\relax}
  >{\raggedright\arraybackslash}p{\dimexpr\linewidth/1000*500-1\tabcolsep\relax}@{}}
\toprule\noalign{}
\endhead
\bottomrule\noalign{}
\endlastfoot
\begin{minipage}[t]{\linewidth}\raggedright
Zatwierdził:

Podpis: \underline{\hspace{3cm}}

Data: \underline{\hspace{4cm}}
\end{minipage} & \begin{minipage}[t]{\linewidth}\raggedright
Stanowisko / Tytuł:

Imię i nazwisko: \underline{\hspace{3cm}}

Funkcja: \underline{\hspace{3cm}}
\end{minipage} \\
\end{longtable}

\phantomsection\label{doc8}
PROC-008

\section{PROCEDURA SZKOLEŃ I WERYFIKACJI WIEDZY Z
CYBERBEZPIECZEŃSTWA}\label{procedura-szkoleux144-i-weryfikacji-wiedzy-z-cyberbezpieczeux144stwa}

\begin{longtable}[]{@{}>{\raggedright\arraybackslash}p{\dimexpr\linewidth/1000*300-2\tabcolsep\relax}>{\raggedright\arraybackslash}p{\dimexpr\linewidth/1000*700-2\tabcolsep\relax}@{}}
\toprule\noalign{}
\endhead
\bottomrule\noalign{}
\endlastfoot
Wersja: & 1.0 \\
Dotyczy: & Wszyscy pracownicy i współpracownicy \\
Odpowiedzialny: & Administrator IT + Właściciel \\
Częstotliwość: & Szkolenie bazowe: onboarding \textbar{} Odświeżające:
co 12 miesięcy \\
\end{longtable}

\subsection{1. Program szkoleń}\label{program-szkoleux144}

\subsubsection{Szkolenie bazowe (nowi
pracownicy)}\label{szkolenie-bazowe-nowi-pracownicy}

\begin{itemize}
\tightlist
\item
  Czas trwania: min. 2 godziny
\item
  Format: prezentacja + omówienie + quiz
\item
  Zawartość: polityki IT, hasła i MFA, phishing i inżynieria społeczna,
  obsługa incydentów, RODO, bezpieczna praca zdalna
\item
  Wymagany wynik quizu: min. \textbf{80\%}
\item
  Powtarzanie przy wyniku poniżej 80\% (max 3 podejścia, potem szkolenie
  indywidualne)
\end{itemize}

\subsubsection{Szkolenie odświeżające (co 12
miesięcy)}\label{szkolenie-odux15bwieux17cajux105ce-co-12-miesiux119cy}

\begin{itemize}
\tightlist
\item
  Czas trwania: min. 1 godzina
\item
  Format: e-learning lub warsztat stacjonarny
\item
  Zawartość: nowe zagrożenia roku, przegląd incydentów z minionego roku,
  nowe procedury
\item
  Wymagany wynik testu: min. \textbf{80\%}
\end{itemize}

\subsection{2. Testy phishingowe}\label{testy-phishingowe}

Firma przeprowadza symulowane ataki phishingowe w celu weryfikacji
świadomości pracowników:

\begin{longtable}[]{@{}>{\raggedright\arraybackslash}p{\dimexpr\linewidth/1000*270-2\tabcolsep\relax}>{\raggedright\arraybackslash}p{\dimexpr\linewidth/1000*430-2\tabcolsep\relax}>{\raggedright\arraybackslash}p{\dimexpr\linewidth/1000*300-2\tabcolsep\relax}@{}}
\toprule\noalign{}
Parametr & Opis & Wartość docelowa \\
\midrule\noalign{}
\endhead
\bottomrule\noalign{}
\endlastfoot
Częstotliwość & Niezapowiedziane testy & Co kwartał (minimum) \\
Narzędzie & Platforma do symulacji phishingu & {[}NAZWA -- np.
GoPhish/Cofense{]} \\
Próg alertu & Odsetek kliknięć w link & Cel: \textless10\%
pracowników \\
Reakcja na kliknięcie & Natychmiastowe szkolenie remediacyjne &
Obowiązkowe w ciągu 7 dni \\
Raportowanie & Anonimowe wyniki zbiorcze & Co kwartał -- kierownik \\
\end{longtable}

\subsection{3. Przegląd Playbooków i
ćwiczenia}\label{przeglux105d-playbookuxf3w-i-ux107wiczenia}

\begin{itemize}
\tightlist
\item
  Pracownicy zapoznają się z Playbookiem Incident Response (PLAY-010) co
  12 miesięcy
\item
  Coroczne ćwiczenie tablicowe (tabletop exercise) -- symulacja
  scenariusza incydentu
\item
  Udokumentowane wyniki ćwiczenia i wnioski (lessons learned)
\end{itemize}

\subsection{4. Rejestr szkoleń}\label{rejestr-szkoleux144}

\begin{longtable}[]{@{}>{\raggedright\arraybackslash}p{\dimexpr\linewidth/1000*220-2\tabcolsep\relax}>{\raggedright\arraybackslash}p{\dimexpr\linewidth/1000*200-2\tabcolsep\relax}>{\raggedright\arraybackslash}p{\dimexpr\linewidth/1000*200-2\tabcolsep\relax}>{\raggedright\arraybackslash}p{\dimexpr\linewidth/1000*200-2\tabcolsep\relax}>{\raggedright\arraybackslash}p{\dimexpr\linewidth/1000*180-2\tabcolsep\relax}@{}}
\toprule\noalign{}
Pracownik & Szkolenie & Data & Wynik testu & Certyfikat \\
\midrule\noalign{}
\endhead
\bottomrule\noalign{}
\endlastfoot
\underline{\hspace{3cm}} & Bazowe & \underline{\hspace{2cm}} &
\underline{\hspace{2cm}}\% & {[} {]} TAK \\
\underline{\hspace{3cm}} & Odświeżające & \underline{\hspace{2cm}} &
\underline{\hspace{2cm}}\% & {[} {]} TAK \\
\underline{\hspace{3cm}} & Phishing test & \underline{\hspace{2cm}}
& Kliknął: {[} {]} TAK / NIE & N/A \\
\underline{\hspace{3cm}} & Bazowe & \underline{\hspace{2cm}} &
\underline{\hspace{2cm}}\% & {[} {]} TAK \\
\end{longtable}

\begin{longtable}[]{@{}
  >{\raggedright\arraybackslash}p{\dimexpr\linewidth/1000*500-1\tabcolsep\relax}
  >{\raggedright\arraybackslash}p{\dimexpr\linewidth/1000*500-1\tabcolsep\relax}@{}}
\toprule\noalign{}
\endhead
\bottomrule\noalign{}
\endlastfoot
\begin{minipage}[t]{\linewidth}\raggedright
Zatwierdził:

Podpis: \underline{\hspace{3cm}}

Data: \underline{\hspace{4cm}}
\end{minipage} & \begin{minipage}[t]{\linewidth}\raggedright
Stanowisko / Tytuł:

Imię i nazwisko: \underline{\hspace{3cm}}

Funkcja: \underline{\hspace{3cm}}
\end{minipage} \\
\end{longtable}

\phantomsection\label{doc9}
POL-009

\section{POLITYKA MONITOROWANIA I AUDYTOWANIA
INFRASTRUKTURY}\label{polityka-monitorowania-i-audytowania-infrastruktury}

\begin{longtable}[]{@{}>{\raggedright\arraybackslash}p{\dimexpr\linewidth/1000*300-2\tabcolsep\relax}>{\raggedright\arraybackslash}p{\dimexpr\linewidth/1000*700-2\tabcolsep\relax}@{}}
\toprule\noalign{}
\endhead
\bottomrule\noalign{}
\endlastfoot
Wersja: & 1.0 \\
Dotyczy: & Wszystkie systemy IT firmy \\
Odpowiedzialny: & Administrator IT \\
Przegląd: & Co 12 miesięcy \\
\end{longtable}

\subsection{1. Logi systemowe -- co
zbieramy}\label{logi-systemowe-co-zbieramy}

\begin{longtable}[]{@{}>{\raggedright\arraybackslash}p{\dimexpr\linewidth/1000*270-2\tabcolsep\relax}>{\raggedright\arraybackslash}p{\dimexpr\linewidth/1000*250-2\tabcolsep\relax}>{\raggedright\arraybackslash}p{\dimexpr\linewidth/1000*250-2\tabcolsep\relax}>{\raggedright\arraybackslash}p{\dimexpr\linewidth/1000*230-2\tabcolsep\relax}@{}}
\toprule\noalign{}
Źródło logów & Rodzaj zdarzeń & Retencja & Szyfrowanie \\
\midrule\noalign{}
\endhead
\bottomrule\noalign{}
\endlastfoot
Serwer Windows/Linux & Logowania, zmiany uprawnień, zdarzenia systemu &
12 miesięcy & {[} {]} TAK \\
Firewall / Router & Połączenia, blokady, anomalie ruchu & 12 miesięcy &
{[} {]} TAK \\
Poczta e-mail & Wysyłane/odebrane, spam, załączniki & 12 miesięcy & {[}
{]} TAK \\
System księgowy & Logowania, zmiany danych, generowanie raportów & 5 lat
& {[} {]} TAK \\
Stacje robocze (EDR) & Procesy, połączenia sieciowe, urządzenia USB & 90
dni & {[} {]} TAK \\
VPN & Sesje, adresy IP, czas trwania & 12 miesięcy & {[} {]} TAK \\
\end{longtable}

\subsection{2. Harmonogram przeglądów i
audytów}\label{harmonogram-przeglux105duxf3w-i-audytuxf3w}

\begin{longtable}[]{@{}>{\raggedright\arraybackslash}p{\dimexpr\linewidth/1000*270-2\tabcolsep\relax}>{\raggedright\arraybackslash}p{\dimexpr\linewidth/1000*250-2\tabcolsep\relax}>{\raggedright\arraybackslash}p{\dimexpr\linewidth/1000*250-2\tabcolsep\relax}>{\raggedright\arraybackslash}p{\dimexpr\linewidth/1000*230-2\tabcolsep\relax}@{}}
\toprule\noalign{}
Rodzaj audytu & Częstotliwość & Odpowiedzialny & Ostatni / Następny \\
\midrule\noalign{}
\endhead
\bottomrule\noalign{}
\endlastfoot
Przegląd logów bezpieczeństwa & Codziennie & Admin IT & \underline{\hspace{2cm}} /
\underline{\hspace{2cm}} \\
Weryfikacja aktywnych kont użytkowników & Co miesiąc & Admin IT & \underline{\hspace{2cm}}
/ \underline{\hspace{2cm}} \\
Audyt uprawnień (równanie z macierzą) & Co kwartał & Admin IT & \underline{\hspace{2cm}} /
\underline{\hspace{2cm}} \\
Skan podatności sieci & Co 6 miesięcy & Admin IT / zewnętrzny & \underline{\hspace{2cm}} /
\underline{\hspace{2cm}} \\
Przegląd konfiguracji firewall & Co 6 miesięcy & Admin IT & \underline{\hspace{2cm}} /
\underline{\hspace{2cm}} \\
Pen-test (test penetracyjny) & Co 12 miesięcy & Zewnętrzny & \underline{\hspace{2cm}} /
\underline{\hspace{2cm}} \\
Audyt zgodności RODO & Co 12 miesięcy & IOD / zewnętrzny & \underline{\hspace{2cm}} /
\underline{\hspace{2cm}} \\
\end{longtable}

\subsection{3. Alerty automatyczne --
triggers}\label{alerty-automatyczne-triggers}

\begin{itemize}
\tightlist
\item
  Logowanie poza godzinami pracy (18:00--6:00 lub weekend) → alert
  natychmiastowy
\item
  5+ nieudanych prób logowania → alert + tymczasowa blokada konta
\item
  Dostęp do plików spoza normalnego zakresu pracy użytkownika
\item
  Wysyłka dużej liczby e-maili przez jedno konto (\textgreater50/godz.)
\item
  Transfer danych \textgreater1 GB w krótkim czasie
\item
  Nowe urządzenie pojawiające się w sieci (rogue device)
\end{itemize}

\subsection{4. Procedura przechowywania i ochrony
logów}\label{procedura-przechowywania-i-ochrony-loguxf3w}

\begin{itemize}
\tightlist
\item
  Logi przechowywane na oddzielnym, zabezpieczonym serwerze lub SIEM
\item
  Dostęp do logów: wyłącznie administrator IT i właściciel firmy
\item
  Logi niemodyfikowalne (write-once lub podpisane cyfrowo)
\item
  Kopia logów offsite lub w chmurze (zaszyfrowana)
\end{itemize}

\begin{longtable}[]{@{}
  >{\raggedright\arraybackslash}p{\dimexpr\linewidth/1000*500-1\tabcolsep\relax}
  >{\raggedright\arraybackslash}p{\dimexpr\linewidth/1000*500-1\tabcolsep\relax}@{}}
\toprule\noalign{}
\endhead
\bottomrule\noalign{}
\endlastfoot
\begin{minipage}[t]{\linewidth}\raggedright
Zatwierdził:

Podpis: \underline{\hspace{3cm}}

Data: \underline{\hspace{4cm}}
\end{minipage} & \begin{minipage}[t]{\linewidth}\raggedright
Stanowisko / Tytuł:

Imię i nazwisko: \underline{\hspace{3cm}}

Funkcja: \underline{\hspace{3cm}}
\end{minipage} \\
\end{longtable}

\phantomsection\label{doc10}
PLAY-010

\section{PLAYBOOK SYTUACJI NIETYPOWYCH -- INCIDENT
RESPONSE}\label{playbook-sytuacji-nietypowych-incident-response}

\begin{longtable}[]{@{}>{\raggedright\arraybackslash}p{\dimexpr\linewidth/1000*300-2\tabcolsep\relax}>{\raggedright\arraybackslash}p{\dimexpr\linewidth/1000*700-2\tabcolsep\relax}@{}}
\toprule\noalign{}
\endhead
\bottomrule\noalign{}
\endlastfoot
Wersja: & 1.0 \\
Dotyczy: & Wszyscy pracownicy \\
Odpowiedzialny: & Właściciel / Administrator IT \\
Aktualizacja: & Po każdym incydencie i co 12 miesięcy \\
\end{longtable}

\textbf{ALARM:} Ten dokument musi być dostępny OFFLINE (wydrukowany) w widocznym
miejscu. W przypadku incydentu sieć może być niedostępna.

\subsection{1. Kontakty alarmowe -- WYPEŁNIĆ I
WYDRUKOWAĆ}\label{kontakty-alarmowe-wypeux142niux107-i-wydrukowaux107}

\begin{longtable}[]{@{}>{\raggedright\arraybackslash}p{\dimexpr\linewidth/1000*270-2\tabcolsep\relax}>{\raggedright\arraybackslash}p{\dimexpr\linewidth/1000*250-2\tabcolsep\relax}>{\raggedright\arraybackslash}p{\dimexpr\linewidth/1000*250-2\tabcolsep\relax}>{\raggedright\arraybackslash}p{\dimexpr\linewidth/1000*230-2\tabcolsep\relax}@{}}
\toprule\noalign{}
Rola & Imię i Nazwisko & Telefon / E-mail & Zakres odpowiedzialności \\
\midrule\noalign{}
\endhead
\bottomrule\noalign{}
\endlastfoot
Właściciel firmy (CISO) & \underline{\hspace{3cm}} &
\underline{\hspace{3cm}} & Główna decyzja, autoryzacja działań \\
Administrator IT & \underline{\hspace{3cm}} &
\underline{\hspace{3cm}} & Techniczne działanie i izolacja \\
Osoba do kontaktu z Policją & \underline{\hspace{3cm}} &
\underline{\hspace{3cm}} & Zgłoszenia prawne, przestępstwa \\
Osoba do kontaktu z mediami & \underline{\hspace{3cm}} &
\underline{\hspace{3cm}} & JEDEN głos firmy, komunikacja zewn. \\
Osoba do kontaktu z UODO & \underline{\hspace{3cm}} &
\underline{\hspace{3cm}} & Naruszenia RODO (72h) \\
Prawnik firmy & \underline{\hspace{3cm}} &
\underline{\hspace{3cm}} & Porady prawne, kontakty z organami \\
Ubezpieczyciel (cyber) & \underline{\hspace{3cm}} &
\underline{\hspace{3cm}} & Zgłoszenie szkody \\
Zewnętrzna firma IT (backup) & \underline{\hspace{3cm}} &
\underline{\hspace{3cm}} & Wsparcie techniczne awaryjne \\
\end{longtable}

Ważne numery: Policja: 997 \textbar{} CERT Polska: cert.pl \textbar{}
UODO: uodo.gov.pl \textbar{} tel. UODO: 22 531 03 00

\subsection{2. Fazy reagowania na
incydent}\label{fazy-reagowania-na-incydent}

\subsubsection{FAZA 1: Identyfikacja (0--30
min)}\label{faza-1-identyfikacja-030-min}

\begin{enumerate}
\tightlist
\item
  Pracownik wykrywa anomalię → \textbf{NATYCHMIASTOWO} informuje
  administratora IT i właściciela
\item
  Administrator IT ocenia powagę incydentu wg klasyfikacji poniżej
\item
  Rejestracja incydentu w Rejestrze Incydentów (godzina, opis, reporter)
\end{enumerate}

\subsubsection{FAZA 2: Izolacja (30--60
min)}\label{faza-2-izolacja-3060-min}

\begin{enumerate}
\tightlist
\item
  Odizolowanie dotkniętych systemów od sieci (jeśli dotyczy) -- kabel
  LAN + WiFi
\item
  Zmiana haseł do zagrożonych kont
\item
  Zabezpieczenie dowodów (logi, kopie plików) -- \textbf{NIE USUWAĆ
  niczego}
\end{enumerate}

\subsubsection{FAZA 3: Ograniczenie i
likwidacja}\label{faza-3-ograniczenie-i-likwidacja}

\begin{enumerate}
\tightlist
\item
  Blokada ruchu sieciowego do/ze źródeł ataku na firewall
\item
  Przywracanie z kopii zapasowej (wg DRP-011)
\item
  Identyfikacja wektora ataku i usunięcie złośliwego kodu
\end{enumerate}

\subsection{3. Matryca powiadamiania}\label{matryca-powiadamiania}

\begin{longtable}[]{@{}>{\raggedright\arraybackslash}p{\dimexpr\linewidth/1000*270-2\tabcolsep\relax}>{\raggedright\arraybackslash}p{\dimexpr\linewidth/1000*250-2\tabcolsep\relax}>{\raggedright\arraybackslash}p{\dimexpr\linewidth/1000*250-2\tabcolsep\relax}>{\raggedright\arraybackslash}p{\dimexpr\linewidth/1000*230-2\tabcolsep\relax}@{}}
\toprule\noalign{}
Kogo powiadomić & Kiedy & Kto powiadamia & Termin \\
\midrule\noalign{}
\endhead
\bottomrule\noalign{}
\endlastfoot
UODO (Urząd Ochrony Danych) & Naruszenie danych osobowych z ryzykiem &
Właściciel / IOD & \textbf{72 godziny} \\
Klienci poszkodowani & Wysokie ryzyko dla ich danych & Właściciel & Bez
zwłoki \\
Policja / CERT Polska & Przestępstwo, atak hakerski, ransomware &
Wyznaczona osoba & Niezwłocznie \\
Media / PR & Jeśli incydent jest/może być publiczny & Właściciel (JEDEN
głos) & Po konsultacji z prawnikiem \\
Ubezpieczyciel & Szkoda objęta polisą cyber & Właściciel & 24--48
godzin \\
\end{longtable}

\subsection{4. Scenariusze -- reakcje}\label{scenariusze-reakcje}

\subsubsection{Scenariusz A: Phishing / kompromitacja
konta}\label{scenariusz-a-phishing-kompromitacja-konta}

\begin{itemize}
\tightlist
\item
  Natychmiast zmień hasło i zresetuj MFA dla konta
\item
  Przegląd aktywności konta (logi) za ostatnie 30 dni
\item
  Sprawdź czy dane były eksportowane lub przesyłane
\item
  Poinformuj pozostałych pracowników (alert phishingowy)
\end{itemize}

\subsubsection{Scenariusz B:
Ransomware}\label{scenariusz-b-ransomware}

\begin{itemize}
\tightlist
\item
  \textbf{NATYCHMIASTOWO} odłącz zainfekowane urządzenie od sieci (kabel
  + WiFi)
\item
  \textbf{NIE PŁACIĆ OKUPU} bez konsultacji z prawnikiem i
  ubezpieczycielem
\item
  Uruchom procedurę DRP-011
\item
  Zgłoś do Policji i CERT Polska (cert.pl)
\item
  Zidentyfikuj wektor ataku i załataj podatność
\end{itemize}

\subsubsection{Scenariusz C: Wyciek danych
klientów}\label{scenariusz-c-wyciek-danych-klientuxf3w}

\begin{itemize}
\tightlist
\item
  Zidentyfikuj jakie dane wyciekły i ile osób dotyczy
\item
  Ocena ryzyka wg RODO -- czy zgłosić do UODO (72h)?
\item
  Powiadom poszkodowanych klientów
\item
  Zabezpiecz dowody (nie usuwaj logów)
\end{itemize}

\subsection{5. Rejestr incydentów
(szablon)}\label{rejestr-incydentuxf3w-szablon}

\begin{longtable}[]{@{}>{\raggedright\arraybackslash}p{\dimexpr\linewidth/1000*220-2\tabcolsep\relax}>{\raggedright\arraybackslash}p{\dimexpr\linewidth/1000*200-2\tabcolsep\relax}>{\raggedright\arraybackslash}p{\dimexpr\linewidth/1000*200-2\tabcolsep\relax}>{\raggedright\arraybackslash}p{\dimexpr\linewidth/1000*200-2\tabcolsep\relax}>{\raggedright\arraybackslash}p{\dimexpr\linewidth/1000*180-2\tabcolsep\relax}@{}}
\toprule\noalign{}
Data/Czas & Opis incydentu & Działania podjęte & Status & Zamknięty \\
\midrule\noalign{}
\endhead
\bottomrule\noalign{}
\endlastfoot
\underline{\hspace{2cm}} &
\underline{\hspace{4cm}} &
\underline{\hspace{4cm}} & {[} {]} Otwarty &
\underline{\hspace{2cm}} \\
\underline{\hspace{2cm}} &
\underline{\hspace{4cm}} &
\underline{\hspace{4cm}} & {[} {]} Otwarty &
\underline{\hspace{2cm}} \\
\end{longtable}

\begin{longtable}[]{@{}
  >{\raggedright\arraybackslash}p{\dimexpr\linewidth/1000*500-1\tabcolsep\relax}
  >{\raggedright\arraybackslash}p{\dimexpr\linewidth/1000*500-1\tabcolsep\relax}@{}}
\toprule\noalign{}
\endhead
\bottomrule\noalign{}
\endlastfoot
\begin{minipage}[t]{\linewidth}\raggedright
Zatwierdził:

Podpis: \underline{\hspace{3cm}}

Data: \underline{\hspace{4cm}}
\end{minipage} & \begin{minipage}[t]{\linewidth}\raggedright
Stanowisko / Tytuł:

Imię i nazwisko: \underline{\hspace{3cm}}

Funkcja: \underline{\hspace{3cm}}
\end{minipage} \\
\end{longtable}

\phantomsection\label{doc11}
DRP-011

\section{PLAN ODTWARZANIA PO KATASTROFIE (DISASTER RECOVERY
PLAN)}\label{plan-odtwarzania-po-katastrofie-disaster-recovery-plan}

\begin{longtable}[]{@{}>{\raggedright\arraybackslash}p{\dimexpr\linewidth/1000*300-2\tabcolsep\relax}>{\raggedright\arraybackslash}p{\dimexpr\linewidth/1000*700-2\tabcolsep\relax}@{}}
\toprule\noalign{}
\endhead
\bottomrule\noalign{}
\endlastfoot
Wersja: & 1.0 \\
Odpowiedzialny: & Właściciel / Administrator IT \\
RTO (Recovery Time Objective): & {[}WYPEŁNIĆ -- np. 4 godziny{]} \\
RPO (Recovery Point Objective): & {[}WYPEŁNIĆ -- np. 24 godziny{]} \\
Testowanie planu: & Co 12 miesięcy (test pełny) + co 3 miesiące (test
częściowy) \\
\end{longtable}

\subsection{1. Klasyfikacja systemów
krytycznych}\label{klasyfikacja-systemuxf3w-krytycznych}

\begin{longtable}[]{@{}>{\raggedright\arraybackslash}p{\dimexpr\linewidth/1000*270-2\tabcolsep\relax}>{\raggedright\arraybackslash}p{\dimexpr\linewidth/1000*250-2\tabcolsep\relax}>{\raggedright\arraybackslash}p{\dimexpr\linewidth/1000*250-2\tabcolsep\relax}>{\raggedright\arraybackslash}p{\dimexpr\linewidth/1000*230-2\tabcolsep\relax}@{}}
\toprule\noalign{}
System & Priorytet & RTO (maks.) & RPO (maks.) \\
\midrule\noalign{}
\endhead
\bottomrule\noalign{}
\endlastfoot
System księgowy (główna aplikacja) & [P1] KRYTYCZNY (P1) & 2 godziny & 4
godziny \\
Baza danych klientów & [P1] KRYTYCZNY (P1) & 2 godziny & 4 godziny \\
Poczta e-mail & [P2] WYSOKI (P2) & 4 godziny & 24 godziny \\
Współdzielone pliki (serwer) & [P2] WYSOKI (P2) & 4 godziny & 24
godziny \\
VPN / Zdalny dostęp & [P3] ŚREDNI (P3) & 8 godzin & 48 godzin \\
Strona internetowa firmy & [P4] NISKI (P4) & 48 godzin & 7 dni \\
\end{longtable}

\subsection{2. Kolejność przywracania
systemów}\label{kolejnoux15bux107-przywracania-systemuxf3w}

\begin{enumerate}
\tightlist
\item
  \textbf{BEZPIECZEŃSTWO} -- czy środowisko jest bezpieczne? (brak
  aktywnego ataku)
\item
  Infrastruktura sieciowa (router, switch, firewall) -- konfiguracja z
  backupu
\item
  Serwer baz danych -- przywracanie z ostatniego czystego backupu
\item
  System księgowy -- weryfikacja integralności danych
\item
  Konta użytkowników i uprawnienia (wg macierzy POL-005)
\item
  Poczta e-mail
\item
  VPN i dostęp zdalny
\item
  Weryfikacja kompletności -- porównanie z RPO
\item
  Informacja dla pracowników o wznowieniu pracy
\item
  Dokumentacja incydentu, wnioski (lessons learned), aktualizacja DRP
\end{enumerate}

\subsection{3. Lokalizacje kopii
zapasowych}\label{lokalizacje-kopii-zapasowych}

\begin{longtable}[]{@{}>{\raggedright\arraybackslash}p{\dimexpr\linewidth/1000*270-2\tabcolsep\relax}>{\raggedright\arraybackslash}p{\dimexpr\linewidth/1000*250-2\tabcolsep\relax}>{\raggedright\arraybackslash}p{\dimexpr\linewidth/1000*250-2\tabcolsep\relax}>{\raggedright\arraybackslash}p{\dimexpr\linewidth/1000*230-2\tabcolsep\relax}@{}}
\toprule\noalign{}
Rodzaj kopii & Lokalizacja & Dostęp (kto / jak) & Przetestowany \\
\midrule\noalign{}
\endhead
\bottomrule\noalign{}
\endlastfoot
Backup lokalny (NAS) & {[}LOKALIZACJA{]} & {[}DANE DOSTĘPU{]} & {[} {]}
\underline{\hspace{2cm}} \\
Backup offsite (fizyczny) & {[}LOKALIZACJA{]} & {[}OSOBA
ODPOWIEDZIALNA{]} & {[} {]} \underline{\hspace{2cm}} \\
Backup chmurowy & {[}PLATFORMA/URL{]} & {[}DANE DOSTĘPU{]} & {[} {]}
\underline{\hspace{2cm}} \\
\end{longtable}

\subsection{4. Procedura testowania DRP}\label{procedura-testowania-drp}

\begin{itemize}
\tightlist
\item
  \textbf{Test częściowy (co kwartał):} przywracanie wybranego systemu
  lub pliku z backupów -- dokumentacja czasu i wyników
\item
  \textbf{Test pełny (co 12 miesięcy):} symulacja pełnej utraty
  środowiska (możliwe środowisko testowe)
\item
  Dokumentacja wyników testu: osiągnięty RTO, RPO, uchybienia
\item
  Aktualizacja DRP na podstawie wyników testów
\end{itemize}

\begin{longtable}[]{@{}
  >{\raggedright\arraybackslash}p{\dimexpr\linewidth/1000*500-1\tabcolsep\relax}
  >{\raggedright\arraybackslash}p{\dimexpr\linewidth/1000*500-1\tabcolsep\relax}@{}}
\toprule\noalign{}
\endhead
\bottomrule\noalign{}
\endlastfoot
\begin{minipage}[t]{\linewidth}\raggedright
Zatwierdził:

Podpis: \underline{\hspace{3cm}}

Data: \underline{\hspace{4cm}}
\end{minipage} & \begin{minipage}[t]{\linewidth}\raggedright
Stanowisko / Tytuł:

Imię i nazwisko: \underline{\hspace{3cm}}

Funkcja: \underline{\hspace{3cm}}
\end{minipage} \\
\end{longtable}

\phantomsection\label{doc12}
PLAN-012

\section{PLAN CIĄGŁEGO DOSKONALENIA
CYBERBEZPIECZEŃSTWA}\label{plan-ciux105gux142ego-doskonalenia-cyberbezpieczeux144stwa}

\begin{longtable}[]{@{}>{\raggedright\arraybackslash}p{\dimexpr\linewidth/1000*300-2\tabcolsep\relax}>{\raggedright\arraybackslash}p{\dimexpr\linewidth/1000*700-2\tabcolsep\relax}@{}}
\toprule\noalign{}
\endhead
\bottomrule\noalign{}
\endlastfoot
Wersja: & 1.0 \\
Odpowiedzialny: & Właściciel / Administrator IT \\
Cykl przeglądu: & Kwartalny przegląd KPI + roczne podsumowanie i
planowanie \\
\end{longtable}

Cyberbezpieczeństwo nie jest jednorazowym projektem, lecz procesem.
Firma zobowiązuje się do regularnego przeglądu, doskonalenia i
rozwijania swoich praktyk bezpieczeństwa.

\subsection{1. Cykl PDCA dla
cyberbezpieczeństwa}\label{cykl-pdca-dla-cyberbezpieczeux144stwa}

\begin{longtable}[]{@{}>{\raggedright\arraybackslash}p{\dimexpr\linewidth/1000*270-2\tabcolsep\relax}>{\raggedright\arraybackslash}p{\dimexpr\linewidth/1000*430-2\tabcolsep\relax}>{\raggedright\arraybackslash}p{\dimexpr\linewidth/1000*300-2\tabcolsep\relax}@{}}
\toprule\noalign{}
Faza & Działanie & Przykłady aktywności \\
\midrule\noalign{}
\endhead
\bottomrule\noalign{}
\endlastfoot
[P] PLAN (Planowanie) & Ocena ryzyka, ustalenie celów & Roczna ocena
ryzyka, ustalenie KPI bezpieczeństwa \\
[D] DO (Realizacja) & Wdrożenie środków & Nowe narzędzia, szkolenia,
zaktualizowane procedury \\
[C] CHECK (Sprawdzenie) & Monitorowanie i audyt & Testy penetracyjne,
przeglądy logów, KPI \\
[A] ACT (Działanie) & Korekty i ulepszenia & Aktualizacja polityk po
incydencie / audycie \\
\end{longtable}

\subsection{2. Roczna ocena ryzyka
(szablon)}\label{roczna-ocena-ryzyka-szablon}

\begin{longtable}[]{@{}>{\raggedright\arraybackslash}p{\dimexpr\linewidth/1000*270-2\tabcolsep\relax}>{\raggedright\arraybackslash}p{\dimexpr\linewidth/1000*250-2\tabcolsep\relax}>{\raggedright\arraybackslash}p{\dimexpr\linewidth/1000*250-2\tabcolsep\relax}>{\raggedright\arraybackslash}p{\dimexpr\linewidth/1000*230-2\tabcolsep\relax}@{}}
\toprule\noalign{}
Ryzyko & Prawdopodobieństwo (1--5) & Wpływ (1--5) & Planowane
działanie \\
\midrule\noalign{}
\endhead
\bottomrule\noalign{}
\endlastfoot
Phishing / kompromitacja konta & \underline{\hspace{2cm}} & \underline{\hspace{2cm}} &
\underline{\hspace{3cm}} \\
Ransomware & \underline{\hspace{2cm}} & \underline{\hspace{2cm}} &
\underline{\hspace{3cm}} \\
Wyciek danych klientów & \underline{\hspace{2cm}} & \underline{\hspace{2cm}} &
\underline{\hspace{3cm}} \\
Utrata backupów & \underline{\hspace{2cm}} & \underline{\hspace{2cm}} &
\underline{\hspace{3cm}} \\
Nieautoryzowany dostęp wewnętrzny & \underline{\hspace{2cm}} & \underline{\hspace{2cm}} &
\underline{\hspace{3cm}} \\
Awaria infrastruktury & \underline{\hspace{2cm}} & \underline{\hspace{2cm}} &
\underline{\hspace{3cm}} \\
Dostawca / procesor zewnętrzny & \underline{\hspace{2cm}} & \underline{\hspace{2cm}} &
\underline{\hspace{3cm}} \\
\end{longtable}

\subsection{3. KPI -- wskaźniki poziomu
bezpieczeństwa}\label{kpi-wskaux17aniki-poziomu-bezpieczeux144stwa}

\begin{longtable}[]{@{}>{\raggedright\arraybackslash}p{\dimexpr\linewidth/1000*270-2\tabcolsep\relax}>{\raggedright\arraybackslash}p{\dimexpr\linewidth/1000*250-2\tabcolsep\relax}>{\raggedright\arraybackslash}p{\dimexpr\linewidth/1000*250-2\tabcolsep\relax}>{\raggedright\arraybackslash}p{\dimexpr\linewidth/1000*230-2\tabcolsep\relax}@{}}
\toprule\noalign{}
Wskaźnik & Cel & Aktualny wynik & Trend \\
\midrule\noalign{}
\endhead
\bottomrule\noalign{}
\endlastfoot
Czas średniego wykrycia incydentu (MTTD) & \textless{} 4 godziny &
\underline{\hspace{2cm}} & {[} {]} ↑ {[} {]} ↓ \\
Czas średniego reagowania (MTTR) & \textless{} 8 godzin &
\underline{\hspace{2cm}} & {[} {]} ↑ {[} {]} ↓ \\
\% pracowników z aktualnym szkoleniem & 100\% & \underline{\hspace{2cm}} &
{[} {]} ↑ {[} {]} ↓ \\
Współcz. kliknięć w phishing test & \textless{} 10\% &
\underline{\hspace{2cm}} & {[} {]} ↑ {[} {]} ↓ \\
Systemy z aktualnym oprogramowaniem & \textgreater{} 95\% &
\underline{\hspace{2cm}} & {[} {]} ↑ {[} {]} ↓ \\
Backup przywracany pomyślnie & 100\% & \underline{\hspace{2cm}} & {[} {]}
↑ {[} {]} ↓ \\
Krytyczne CVE bez łatki (\textgreater30 dni) & 0 &
\underline{\hspace{2cm}} & {[} {]} ↑ {[} {]} ↓ \\
\end{longtable}

\subsection{4. Harmonogram działań
doskonalenia}\label{harmonogram-dziaux142aux144-doskonalenia}

\begin{longtable}[]{@{}>{\raggedright\arraybackslash}p{\dimexpr\linewidth/1000*270-2\tabcolsep\relax}>{\raggedright\arraybackslash}p{\dimexpr\linewidth/1000*250-2\tabcolsep\relax}>{\raggedright\arraybackslash}p{\dimexpr\linewidth/1000*250-2\tabcolsep\relax}>{\raggedright\arraybackslash}p{\dimexpr\linewidth/1000*230-2\tabcolsep\relax}@{}}
\toprule\noalign{}
Działanie & Termin & Odpowiedzialny & Status \\
\midrule\noalign{}
\endhead
\bottomrule\noalign{}
\endlastfoot
Roczna ocena ryzyka & Styczeń & Właściciel & {[} {]}
\underline{\hspace{2cm}} \\
Przegląd wszystkich polityk i procedur & Styczeń & Admin IT & {[} {]}
\underline{\hspace{2cm}} \\
Test penetracyjny (zewnętrzny) & Q2 & Zewnętrzny dostawca & {[} {]}
\underline{\hspace{2cm}} \\
Test DRP (pełny) & Q3 & Admin IT & {[} {]} \underline{\hspace{2cm}} \\
Szkolenia odświeżające & Co 12 miesięcy & Admin IT & {[} {]}
\underline{\hspace{2cm}} \\
Przegląd i aktualizacja playbooka & Po każdym incydencie & Właściciel &
{[} {]} \underline{\hspace{2cm}} \\
Audyt zgodności RODO & Q4 & IOD / zewnętrzny & {[} {]}
\underline{\hspace{2cm}} \\
\end{longtable}

\begin{longtable}[]{@{}
  >{\raggedright\arraybackslash}p{\dimexpr\linewidth/1000*500-1\tabcolsep\relax}
  >{\raggedright\arraybackslash}p{\dimexpr\linewidth/1000*500-1\tabcolsep\relax}@{}}
\toprule\noalign{}
\endhead
\bottomrule\noalign{}
\endlastfoot
\begin{minipage}[t]{\linewidth}\raggedright
Zatwierdził:

Podpis: \underline{\hspace{3cm}}

Data: \underline{\hspace{4cm}}
\end{minipage} & \begin{minipage}[t]{\linewidth}\raggedright
Stanowisko / Tytuł:

Imię i nazwisko: \underline{\hspace{3cm}}

Funkcja: \underline{\hspace{3cm}}
\end{minipage} \\
\end{longtable}

{[}NAZWA FIRMY KSIĘGOWEJ{]} • Pakiet Dokumentacji Cyberbezpieczeństwa
v1.0 • POUFNE

\end{document}